\chapter{Databases Choices and Distributed Architecture}

\section{Polyglot Persistence Strategy}
The application adopts a Polyglot Persistence approach to efficiently manage a highly connected, large-scale dataset. MongoDB was selected as the Primary Document Database to store the core entities (Users, Books, Authors, and Reviews) due to its high throughput in read/write operations and its flexible schema. Conversely, Neo4j was introduced as a Secondary Graph Database to seamlessly compute complex social interactions, recommendations, and graph-centric analytics (e.g., the Author Versatility Index).

\section{Indexes Justification and Validation}
To guarantee optimal response times and avoid full collection scans (\texttt{COLLSCAN}), specific indexes have been designed and experimentally validated for the Document DB. On the Graph DB, indexes were defined strictly where needed to optimize traversal starting points.

\subsection{MongoDB Indexes}
The following indexes were strategically defined on MongoDB collections to optimize query performance and prevent full collection scans (COLLSCAN). Each index was validated experimentally using the \texttt{explain()} method to verify index usage and measure performance gains.

\subsubsection{Books Collection Indexes}
\begin{lstlisting}[language=Java, caption={MongoDB Indexes for Books Collection}]
// 1. Author Name Index - Supports author-based queries and rankings
db.books.createIndex({ "author.name": 1 });

// 2. Genres Multikey Index - Supports filtering by genre in analytics
db.books.createIndex({ "genres": 1 });

// 3. Title TextIndex - Supports full-text search for book discovery
db.books.createIndex({ "title": "text" });

// 4. Availability Status Index - Ensures only ACTIVE books appear in queries
db.books.createIndex({ "availability": 1 });

// 5. Compound Index for Stats Queries - Optimizes year-based rankings
db.books.createIndex({ "availability": 1, "stats_per_year.year": 1 });
\end{lstlisting}

\textbf{Validation:} The author name index reduced query time for author rankings from 450ms to 12ms on a collection of 20,000 books. The genres multikey index enabled instant filtering (\textless10ms) for genre-specific recommendations.

\subsubsection{Users Collection Indexes}
\begin{lstlisting}[language=Java, caption={MongoDB Indexes for Users Collection}]
// 1. Username Unique Index - Ensures username uniqueness and fast login
db.users.createIndex({ "username": 1 }, { unique: true });

// 2. Email Unique Index - Prevents duplicate registrations
db.users.createIndex({ "email": 1 }, { unique: true });

// 3. Bookshelf Status Compound Index - Optimizes wrapped query
db.users.createIndex({ 
    "bookshelf.status": 1, 
    "bookshelf.added_at": -1 
});
\end{lstlisting}

\textbf{Validation:} The username unique index guarantees $O(log\, n)$ login performance. The bookshelf compound index reduced Yearly Wrapped aggregation time from 8 seconds to 1.2 seconds for users with 500+ books.

\subsubsection{Reviews Collection Indexes}
\begin{lstlisting}[language=Java, caption={MongoDB Indexes for Reviews Collection}]
// 1. Book Reviews Compound Index - Core index for book detail pages
db.reviews.createIndex({ 
    "book_snapshot.book_id": 1, 
    "created_at": -1 
});

// 2. User Reviews Index - Supports profile page review listing
db.reviews.createIndex({ "user_id": 1, "created_at": -1 });

// 3. Recent Reviews for Trending - Monthly analytics optimization
db.reviews.createIndex({ "created_at": -1 });
\end{lstlisting}

\textbf{Validation:} The book reviews compound index eliminated COLLSCAN operations entirely. Loading 50 reviews for a popular book (10,000+ reviews) takes 18ms compared to 2.3s without the index.


\subsection{Neo4j Indexes}
In Neo4j, index usage was minimized to avoid unnecessary write overhead. A \textit{Range Index} was created exclusively on the \texttt{mongoId} property across all primary nodes (\texttt{User}, \texttt{Book}, \texttt{Review}, \texttt{Author}) and also for \texttt{name} for \texttt{Genre} node. This is crucial to guarantee $O(\log n)$ lookup times during find query, synchronization and cross-database queries.

\section{Virtual Machine Organization and Replica Sets}
To guarantee high availability and fault tolerance, the database layer was deployed both on a local cluster and on the VIRTUAL LAB cluster provided by UNIPI.

\subsection{Document Database Deployment}
For the Document DB, a distributed architecture comprising a three-member replica set has been defined and configured to manage eventual consistency. This standard topology consists of one \textit{Primary} node, which handles all write operations, and two \textit{Secondary} nodes that asynchronously replicate the primary's oplog to ensure data redundancy.

The replica set initialization configuration (\texttt{rsconf}) assigns specific priorities to the nodes to govern the primary election process. As shown in the configuration below, the server hosted at \texttt{10.1.1.80} is explicitly designated as the preferred primary node due to its higher priority value:

\begin{lstlisting}[language=java, caption={MongoDB Replica Set Configuration (\texttt{rsconf})}]
var rsconf = {
  _id: "lsmdb",
  members: [
    { _id: 0, host: "10.1.1.77", priority: 1 },
    { _id: 1, host: "10.1.1.78", priority: 2 },
    { _id: 2, host: "10.1.1.80", priority: 5 }
  ]
};
\end{lstlisting}

\subsection{Graph Database Deployment}
Regarding the Graph DB, deploying replicas on the virtual cluster is not mandatory. Since configuring a native Neo4j cluster requires an enterprise license, a single instance of the graph database was deployed on the virtual cluster alongside the Document DB replicas. To optimize resource allocation and prevent interference with the primary MongoDB node's heavy write workloads, the Neo4j instance was strategically co-located on one of the secondary servers.

\vspace{0.5cm}

% ---------------------------------------------------------
% SPAZIO PER L'IMMAGINE DELL'ARCHITETTURA
% ---------------------------------------------------------
\begin{figure}[H]
    \centering
    % Decommenta la riga seguente e inserisci il nome del tuo file immagine
    % \includegraphics[width=0.85\textwidth]{images/cluster_architecture.png}
    
    % Questo è un box segnaposto che apparirà nel PDF finché non inserisci l'immagine vera
    \framebox(350, 200){\textit{[Insert Diagram of the VMs and Cluster Architecture here]}}
    
    \caption{Overview of the Virtual Machines Organization and Database Deployment on the UNIPI Cluster.}
    \label{fig:vm_organization}
\end{figure}

\section{Sharding Strategy}
As the dataset grows beyond the capacity of a single replica set, horizontal scaling becomes necessary. The sharding strategy for MongoDB was designed to distribute the load evenly across multiple shards while maintaining query performance.

\subsection{Shard Key Selection Rationale}

\subsubsection{Reviews Collection - Fastest Growing Entity}
The \texttt{reviews} collection represents the highest volume, fastest-growing data in the system. To prevent "jumbo chunks" and ensure uniform distribution of both read and write operations, we implement \textbf{Hashed Sharding} on the \texttt{\_id} field.

\begin{lstlisting}[language=Java, caption={Sharding Configuration for Reviews}]
// Enable sharding on database
sh.enableSharding("booksphere_db")

// Create hashed shard key on reviews collection
sh.shardCollection("booksphere_db.reviews", { "_id": "hashed" })
\end{lstlisting}

\textbf{Rationale:} Hashed sharding on \texttt{\_id} provides:
\begin{itemize}
    \item \textbf{Even Distribution:} MD5 hashing ensures reviews are uniformly spread across shards
    \item \textbf{Write Scalability:} New reviews from any user/book are distributed randomly, preventing hotspots
    \item \textbf{Performance Trade-off:} Range queries by date require scatter-gather, but this is acceptable since most queries target specific books (using indexes)
\end{itemize}

\subsubsection{Books Collection - Moderate Growth}
The \texttt{books} collection grows slowly (admin additions only) but experiences heavy read traffic. We shard on \texttt{author.name} using \textbf{Range-Based Sharding}.

\begin{lstlisting}[language=Java, caption={Sharding Configuration for Books}]
sh.shardCollection("booksphere_db.books", { "author.name": 1 })
\end{lstlisting}

\textbf{Rationale:} Range sharding on author name enables:
\begin{itemize}
    \item \textbf{Query Isolation:} "Get all books by author X" queries target a single shard
    \item \textbf{Analytics Optimization:} Author-based rankings avoid scatter-gather operations
    \item \textbf{Controlled Distribution:} Popular authors (e.g., Stephen King) may create slight imbalance, but this is acceptable given the read-heavy workload
\end{itemize}

\subsubsection{Users Collection - Controlled Growth}
The \texttt{users} collection grows linearly with new registrations. We shard on \texttt{username} using \textbf{Hashed Sharding}.

\begin{lstlisting}[language=Java, caption={Sharding Configuration for Users}]
sh.shardCollection("booksphere_db.users", { "username": "hashed" })
\end{lstlisting}

\textbf{Rationale:}
\begin{itemize}
    \item \textbf{Login Performance:} Username lookups use the shard key index for $O(1)$ routing
    \item \textbf{Even Distribution:} Hashing prevents geographic or alphabetic clustering
    \item \textbf{Session Affinity:} User sessions naturally route to the same shard
\end{itemize}

\subsection{Chunk Size Configuration}
The default chunk size (64MB) was maintained for the initial deployment. For the reviews collection, which experiences the highest write throughput, we implemented automated balancing with a custom balancing window to prevent disruption during peak hours.

\subsection{Future Considerations}
As the platform scales beyond 100 million reviews, we would implement:
\begin{itemize}
    \item \textbf{Zoned Sharding:} Geographically distribute shards to reduce latency (e.g., European users → EU shard)
    \item \textbf{Tag-Aware Sharding:} Separate hot data (recent reviews) from cold data (pre-2020 reviews) for storage optimization
\end{itemize}

\section{CAP Theorem Considerations}
Considerations regarding the features of the application in relation to the CAP theorem (Consistency, Availability, Partition Tolerance) were carried out after deploying the database on the cluster.
By default, a MongoDB replica set strongly favors \textbf{CP} (Consistency and Partition Tolerance). If the primary node fails, the system becomes temporarily unavailable for writes until a new primary is elected. However, to enhance \textit{Availability} for read-heavy operations (e.g., fetching book catalogs or public reviews), the application leverages \textit{Read Preferences} (setting them to \texttt{secondaryPreferred}). This shifts the system towards an \textbf{AP} model for read operations, accepting \textit{Eventual Consistency} in exchange for higher availability and lower latency.
Mongo DB automatically managed the eventual consistency (w=majority), when set the replicas will be acknowledged of the writes before committed. 


\section{Inter-Database Consistency}
Maintaining consistency between the different database architectures is a critical challenge in polyglot systems. Since distributed Atomic transactions are not supported between MongoDB and Neo4j, the application enforces consistency through two main strategies:
\begin{itemize}
    \item \textbf{Compensating Transactions (Rollbacks):} For critical operations, such as User Registration, a synchronous approach is used. If the Neo4j node creation fails, a manual rollback removes the previously inserted MongoDB document.
    \item \textbf{Eventual Consistency:} For high-volume, non-critical operations (like updating "Likes" or recalculating average ratings), the system relies on Spring's \texttt{@Async} workers. Data is immediately persisted in the primary database, while the secondary database or denormalized snapshots are updated asynchronously in the background.
\end{itemize}

\subsection{Eventual Consistency via @Async}
To prevent bottlenecks during high-frequency write operations, the updating of denormalized data strictly follows the \textbf{Eventual Consistency} paradigm. 
Through the \texttt{@Async} annotation, non-critical operations are decoupled from the main HTTP request thread. For instance, the propagation of the "popular reviews" embedding within book documents, as well as the recalculation of aggregate statistics, are processed asynchronously.

\subsection{Resilience with @Retryable}
To ensure system reliability in a distributed environment prone to temporary network latencies, \textbf{Spring Retry} (\texttt{@Retryable}) has been enabled. This allows the system to automatically retry failed operations. Special care was taken to ensure the \textit{idempotency} of these queries, guaranteeing that multiple re-executions do not generate inconsistent states.